\documentclass{article}
\begin{document}
\section{Bubbles}
Throughout this section all products of two-time functions are pointwise in the contour times,
so that the polarization bubble is a product of two Green's functions taken at the same pair of
arguments. This is the structure that appears in the lowest-order density response, in the
second-order self-energy, and in the current-current correlation function, and in each case the
real-time components follow from the analytic continuation of the contour expression. We now
discuss the physical content of these components in some detail. The lesser component measures
the occupied part of the spectrum and enters the particle current; the greater component measures
the empty part and enters the hole current; their difference is the spectral function of the
composite object. In equilibrium these are related by the fluctuation-dissipation theorem, and
out of equilibrium the lesser and greater functions carry independent information about the
distribution. For a steady state they depend only on the time difference, and their Fourier
transforms give the emission and absorption spectra. With these remarks in mind we turn to the
explicit rules. For $C = AB$ the Langreth rules give
\begin{equation}
C^{<} = A^{<} B^{<} .
\label{eq:pl}
\end{equation}
\end{document}
